%------------------------------------------------------------------------------%
\documentclass[a4paper,12pt,fleqn]{article}

\usepackage{integracao_series}

\ShowAnswers

%------------------------------------------------------------------------------%
\begin{document}

\makeHeader{Integração e Séries}{Prova 3 -- Turma 3}{08/06/2026}

% 01
%------------------------------------------------------------------------------%
\exercise{25\%}
Calcule a soma da série
\(
  \sum_{n=1}^{\infty} \frac{3}{2^n}
\)

\begin{answer}
  Esta é uma série geométrica
  \[
    S
    = \sum_{n=1}^{\infty} \frac{3}{2^n}
    = \sum_{n=1}^{\infty} \frac{3}{2} \; \left(\frac{1}{2}\right)^{n-1}
  \]
  Como \(a = \frac{3}{2}\), \(r = \frac{1}{2}\) e \(\abs{r} < 1\)
  \[
    S
     = \frac{a}{1-r}
     = \frac{\sfrac{3}{2}}{1-\sfrac{1}{2}}
     = \frac{3}{\;2\; \dfrac{2-1}{2}\;}
     = \frac{3}{\;2\; \dfrac{1}{2}\;}
     = 3
  \]
\end{answer}

% 02
%------------------------------------------------------------------------------%
\exercise{25\%}
Calcule o limite da sequência de somas parciais da série
\(
  \sum_{n=1}^{\infty} \frac{2}{2n-1} - \frac{2}{2n+1}
\)

\begin{answer}
  A soma parcial de ordem \(N\) é
  \[
    S_N = \sum_{n=1}^{N} \frac{2}{2n-1} - \frac{2}{2n+1}
  \]

  Escrevendo os primeiros termos,
  \[
    S_N
    = \left(       2        - \frac{2}{3}    \right)
    + \left( \frac{2}{3}    - \frac{2}{5}    \right)
    + \left( \frac{2}{5}    - \frac{2}{7}    \right) + \cdots
    + \left( \frac{2}{2N-1} - \frac{2}{2N+1} \right)
  \]
  notamos que
  \[
    S_N = 2 - \frac{2}{2N+1}
  \]
  portanto
  \[
    \lim_{N\to\infty} S_N
    = \lim_{N\to\infty} \left(2 - \frac{2}{2N+1}\right)
    = 2 - \lim_{N\to\infty} \frac{2}{2N+1}
    = 2 - 0
    = 2
  \]
\end{answer}

% 03
%------------------------------------------------------------------------------%
\exercise{25\%}
Utilize o teste da integral para verificar a convergência da série
\(
  \sum_{n=1}^{\infty} \left(\frac{1}{n}\right)^3
\)\\
Verifique todas as condições do teste

\begin{answer}
  Escolhemos a função \(f(x) = \frac{1}{x^3}\), pois
  \[
    a_n = f(n) = \frac{1}{n^3} = \left(\frac{1}{n}\right)^3
  \]
  Para $x\in[1, \infty)$ a \textbf{função} $f(x)$ é decrescente, contínua e positiva.

  Calculando a integral imprópria
  \begin{align*}
    \int_{1}^{\infty} \frac{1}{x^3}\,dx
    & = \lim_{b\to\infty} \int_{1}^{b} x^{-3}\,dx                                       \\[2mm]
    & = \lim_{b\to\infty} \left[ \;\frac{\,x^{1-3}\,}{\,1-3\,}\evalat{1}{b}\;\, \right] \\[2mm]
    & = \lim_{b\to\infty} \left[ \frac{b^{-2}}{-2} - \frac{b^0}{-2} \right]             \\[2mm]
    & = 0 + \frac{1}{2}
  \end{align*}
  Pelo \textbf{teste da integral}, como a integral imprópria converge a série também converge
\end{answer}

% 04
%------------------------------------------------------------------------------%
\exercise{25\%}
Verifique a convergência da série
\(
  \sum_{n=1}^{\infty}\frac{n!}{3^n}
\)
\begin{answer}
  \textbf{Solução 1: pelo teste da razão}

  Como
  \[
    a_n=\frac{n!}{3^n}
  \]
  temos
  \[
    a_{n+1}
    = \frac{(n+1)!}{3^{n+1}}
  \]

  Como \(a_n>0\) para todo \(n\), podemos aplicar o teste da razão
  \[
    \frac{a_{n+1}}{a_n}
    = \frac{\dfrac{(n+1)!}{3^{n+1}}}{\dfrac{n!}{3^n}}
    = \frac{(n+1)!}{3^{n+1}}\cdot\frac{3^n}{n!}
    = \frac{(n+1)\,n!\,3^n}{3\cdot 3^n\,n!}
    = \frac{n+1}{3}
  \]
  Calculando $\rho$
  \[
    \rho
    = \lim_{n\to\infty} \frac{a_{n+1}}{a_n}
    = \lim_{n\to\infty} \frac{n+1}{3}
    = \infty
  \]
  O \textbf{teste da razão} garante que a série diverge

  \bigskip
  \bigskip
  \textbf{Solução 2: pelo teste da divergência}
  \[
    \lim_{n\to\infty} \frac{n!}{3^n}
    = \lim_{n\to\infty} \frac{n(n-1)(n-2) \cdots 4 \cdot 3 \cdot 2 \cdot 1 }
                             {3 \cdot 3 \cdot 3 \cdots 3 \cdot 3 \cdot 3 \cdot 3}
   = \lim_{n\to\infty}
    \left(
      \frac{n}{3}
      \frac{n-1}{3}
      \frac{n-2}{3}
      \cdots
      \frac{4}{3}
      \frac{3}{3}
      \frac{2}{3}
      \frac{1}{3}
    \right)
    = \infty
  \]
  Portanto a série é divergente pelo \textbf{teste da divergência}

\end{answer}

%------------------------------------------------------------------------------%
\end{document}
%------------------------------------------------------------------------------%
